\documentclass[12pt,a4paper]{article}
\usepackage[top=2.5cm,bottom=2.5cm,left=2.5cm,right=2.5cm,headheight=15pt]{geometry}
\usepackage{amsmath,amssymb}
\usepackage{tikz}
\usetikzlibrary{shapes,arrows,arrows.meta,positioning,calc}
\usepackage{booktabs,array}
\usepackage{xcolor}
\usepackage{enumitem}
\usepackage{fancyhdr}
\usepackage{titlesec}
\usepackage{mdframed}
\usepackage{tcolorbox}
\tcbuselibrary{skins,breakable}
\usepackage{hyperref}

%% ---- Colors ----
\definecolor{folblue}{RGB}{0,82,165}
\definecolor{lightblue}{RGB}{220,235,255}
\definecolor{darkorange}{RGB}{180,70,0}
\definecolor{lightorange}{RGB}{255,235,210}
\definecolor{darkgreen}{RGB}{20,110,20}
\definecolor{lightgreen}{RGB}{200,240,205}
\definecolor{darkred}{RGB}{150,20,20}
\definecolor{lightred}{RGB}{255,210,210}
\definecolor{lightgray}{RGB}{242,242,242}

%% ---- tcolorbox styles ----
\newtcolorbox{qbox}[2][]{enhanced,breakable,
  colback=lightorange,colframe=darkorange,arc=4pt,
  title={\textbf{#2}},fonttitle=\bfseries\color{white},
  attach boxed title to top left={yshift=-2mm,xshift=4mm},
  boxed title style={colback=darkorange,arc=3pt},#1}

\newtcolorbox{solbox}[2][]{enhanced,breakable,
  colback=lightblue,colframe=folblue,arc=4pt,
  title={\textbf{#2}},fonttitle=\bfseries\color{white},
  attach boxed title to top left={yshift=-2mm,xshift=4mm},
  boxed title style={colback=folblue,arc=3pt},#1}

\newtcolorbox{keybox}[2][]{enhanced,breakable,
  colback=lightgreen,colframe=darkgreen,arc=4pt,
  title={\textbf{#2}},fonttitle=\bfseries\color{white},
  attach boxed title to top left={yshift=-2mm,xshift=4mm},
  boxed title style={colback=darkgreen,arc=3pt},#1}

\newtcolorbox{algobox}[2][]{enhanced,breakable,
  colback=lightgray,colframe=gray!60,arc=4pt,
  title={\textbf{#2}},fonttitle=\bfseries,
  attach boxed title to top left={yshift=-2mm,xshift=4mm},
  boxed title style={colback=gray!60,arc=3pt},#1}

\newtcolorbox{warnbox}[2][]{enhanced,breakable,
  colback=lightred,colframe=darkred,arc=4pt,
  title={\textbf{#2}},fonttitle=\bfseries\color{white},
  attach boxed title to top left={yshift=-2mm,xshift=4mm},
  boxed title style={colback=darkred,arc=3pt},#1}

%% ---- Page style ----
\pagestyle{fancy}\fancyhf{}
\lhead{\color{folblue}\textbf{Markov Chains --- CSE 717 InfoSec}}
\rhead{\color{folblue}Exam: 17 Jun 2026}
\cfoot{--- \thepage\ ---}
\renewcommand{\headrulewidth}{0.4pt}

%% ---- Section style ----
\titleformat{\section}{\large\bfseries\color{folblue}}{}{0em}{}[\titlerule]
\titleformat{\subsection}{\normalsize\bfseries\color{darkorange}}{}{0em}{}

\begin{document}

\begin{center}
{\LARGE\bfseries\color{folblue} CSE 717 --- Information Security}\\[4pt]
{\large\bfseries Topic 6: Markov Chains --- Steady-State Probability, 2-State Security Model}\\[2pt]
{\large Complete Past Paper Solutions: 2020--2024}\\[6pt]
{\normalsize Source: Lc\#4A (Markov Chain theory) + Lc\#4B (Assignment numericals)}\\[2pt]
{\normalsize\color{darkred}\textbf{2/5 years (2023, 2024 --- both most recent), 3--4 marks. Strong recency signal --- near-certain repeat.}}
\end{center}

\tableofcontents

%% ==========================================
\section*{Quick Reference --- Steady-State Method (CHEAT SHEET)}
%% ==========================================

\begin{warnbox}{2020--2022: no Markov Chain question found}
This topic first appeared in 2023 and 2024 (both most recent years) as a 2-state security model: \textbf{Secure/Compromised} or \textbf{Secure/Insecure}. The method is identical both years --- only the transition probabilities change. Master one worked example cold.
\end{warnbox}

\begin{algobox}{Key Definitions (Lc\#4A)}
\begin{itemize}[noitemsep]
  \item \textbf{Markov chain:} a model where the next state depends \emph{only} on the current state (memoryless / Markov property).
  \item \textbf{State:} a possible situation of the system (e.g.\ Secure, Compromised, Sunny, Rainy).
  \item \textbf{Transition probability:} the probability of moving from one state to another in one step. All outgoing probabilities from a state must sum to 1.
  \item \textbf{Transition matrix $T$:} $k \times k$ matrix for $k$ states; entry $T_{ij}$ = probability of going from state $i$ to state $j$. Each \emph{row} sums to 1.
  \item \textbf{State vector $x^{(n)}$:} a row vector of probabilities of being in each state at step $n$. Updated as $x^{(n)} = x^{(n-1)} T$.
  \item \textbf{Steady-state vector $V$:} the long-run probability distribution, satisfying $VT = V$ (does not change after further steps). Also written $V(T - I) = 0$ with the constraint $\sum V_i = 1$.
\end{itemize}
\end{algobox}

\begin{algobox}{5-Step Algorithm to Find Steady-State (Lc\#4A standard method)}
\begin{enumerate}[noitemsep]
  \item \textbf{Set up transition matrix $T$} from the given probabilities (rows = ``from'', columns = ``to'').
  \item \textbf{Write $V(T-I) = 0$}: subtract the identity matrix $I$ from $T$, multiply row vector $V = [S \;\; C]$ on the left, set equal to zero row vector.
  \item \textbf{Extract equations}: one equation per column (they are linearly dependent --- any one can be dropped). This gives a relation between $S$ and $C$.
  \item \textbf{Use the normalization constraint} $S + C = 1$ (probabilities must sum to 1) together with the relation from Step 3.
  \item \textbf{Solve for $S$ and $C$}, state the steady-state vector $V = [S \;\; C]$, and \textbf{interpret} (which state dominates long-run?).
\end{enumerate}
\end{algobox}

\begin{keybox}{Worked example from slides (Lc\#4A) --- Sunny/Rainy weather}
\[
T = \begin{pmatrix} 0.9 & 0.1 \\ 0.5 & 0.5 \end{pmatrix}
\quad \text{(Sunny$\to$Sunny=0.9, Rainy$\to$Sunny=0.5)}
\]
$T - I = \begin{pmatrix} -0.1 & 0.1 \\ 0.5 & -0.5 \end{pmatrix}$

$[S\;\;R]\begin{pmatrix}-0.1 & 0.1 \\ 0.5 & -0.5\end{pmatrix} = [0\;\;0]$
$\;\Rightarrow\; -0.1S + 0.5R = 0 \;\Rightarrow\; S = 5R$

With $S + R = 1$: $5R + R = 1 \Rightarrow R = 1/6 \approx 0.167$, $S = 5/6 \approx 0.833$.

$\boxed{V = [0.833\;\;0.167]}$ --- long-run: 83.3\% sunny days.
\end{keybox}

%% ==========================================
\section{2024 --- Q1(c): Secure/Compromised 2-state system --- find steady-state, interpret (3)}
%% ==========================================

\begin{qbox}{Question --- 2024 Q1(c) \hfill [3 marks]}
An information system can be in one of two states: \textbf{Secure (S)} or \textbf{Compromised (C)}.
\begin{itemize}[noitemsep]
  \item If the system is \textbf{Secure} during a security assessment cycle (active patching, effective controls), it remains Secure with probability \textbf{0.4}. Otherwise it transitions to Compromised.
  \item If the system is \textbf{Compromised} during an assessment cycle (failed controls, unpatched vulnerabilities), incident-response actions return it to Secure with probability \textbf{0.3}. Otherwise it remains Compromised.
\end{itemize}
Determine the steady-state probabilities of Secure (S) and Compromised (C) states, and evaluate whether the system is more likely to be secure or compromised in the long run.
\end{qbox}

\begin{solbox}{Answer --- 2024 Q1(c)}

\textbf{Step 1 --- Transition Matrix.}

From the given probabilities:
\begin{itemize}[noitemsep]
  \item Secure $\to$ Secure: $0.4$;\quad Secure $\to$ Compromised: $1 - 0.4 = 0.6$
  \item Compromised $\to$ Secure: $0.3$;\quad Compromised $\to$ Compromised: $1 - 0.3 = 0.7$
\end{itemize}
\[
T = \begin{pmatrix} 0.4 & 0.6 \\ 0.3 & 0.7 \end{pmatrix}
\quad\text{(rows = ``from'' state, columns = ``to'' state)}
\]

\textbf{Step 2 --- Set up steady-state equation $V(T - I) = 0$.}
\[
T - I = \begin{pmatrix} 0.4 & 0.6 \\ 0.3 & 0.7 \end{pmatrix} - \begin{pmatrix} 1 & 0 \\ 0 & 1 \end{pmatrix} = \begin{pmatrix} -0.6 & 0.6 \\ 0.3 & -0.3 \end{pmatrix}
\]
\[
[S \;\; C]\begin{pmatrix} -0.6 & 0.6 \\ 0.3 & -0.3 \end{pmatrix} = [0 \;\; 0]
\]

\textbf{Step 3 --- Extract equation from column 1.}
\[
-0.6S + 0.3C = 0 \;\Rightarrow\; C = 2S \quad \cdots (1)
\]
(Column 2 gives $0.6S - 0.3C = 0$ --- the same equation, linearly dependent. Discard it.)

\textbf{Step 4 --- Apply normalization constraint $S + C = 1$.}
\[
S + C = 1 \quad \cdots (2)
\]

\textbf{Step 5 --- Solve.}

Substituting (1) into (2):
\[
S + 2S = 1 \;\Rightarrow\; 3S = 1 \;\Rightarrow\; \boxed{S = \tfrac{1}{3} \approx 0.333}
\]
\[
C = 2S = \tfrac{2}{3} \;\Rightarrow\; \boxed{C \approx 0.667}
\]

\textbf{Steady-state vector:}
\[
V = [S \;\; C] = [0.333 \;\; 0.667]
\]

\textbf{Interpretation:} In the long run, the system is in the \textbf{Secure} state only \textbf{33.3\%} of the time and in the \textbf{Compromised} state \textbf{66.7\%} of the time.

\[
\boxed{\text{The system is more likely COMPROMISED in the long run (66.7\% vs 33.3\%).}}
\]

\textit{Intuition: although incident-response recovers from Compromised state 30\% of the time, the patching mechanism only keeps the system Secure 40\% of the time --- so the system drifts towards compromise.}
\end{solbox}

%% ==========================================
\section{2023 --- Section B Q6(c): Secure/Insecure network with patch-cycle transitions --- long-term behaviour (4)}
%% ==========================================

\begin{qbox}{Question --- 2023 Section B Q6(c) \hfill [4 marks]}
Consider a network system that can either be in a \textbf{secure state} or an \textbf{insecure state}. The network system periodically undergoes patches to maintain security. However, there is a risk that the system may be vulnerable to attacks at any given time.
\begin{itemize}[noitemsep]
  \item[(i)] When the system remains secure because the patching mechanism is effective and there are no new vulnerabilities, there is an \textbf{80\% chance} that the system remains secure after a security check.
  \item[(ii)] When no patch is applied or the patching mechanism fails, the system remains insecure and there is a \textbf{60\% chance} that the system remains secure after a security check.
\end{itemize}
Determine the long-term behaviour of this network system, whether it is secure or insecure.
\end{qbox}

\begin{solbox}{Answer --- 2023 Section B Q6(c)}

\textbf{Step 1 --- Transition Matrix.}

From the given probabilities:
\begin{itemize}[noitemsep]
  \item Secure $\to$ Secure: $0.8$;\quad Secure $\to$ Insecure: $1 - 0.8 = 0.2$
  \item Insecure $\to$ Secure: $0.6$;\quad Insecure $\to$ Insecure: $1 - 0.6 = 0.4$
\end{itemize}
\[
T = \begin{pmatrix} 0.8 & 0.2 \\ 0.6 & 0.4 \end{pmatrix}
\quad\text{(rows = ``from'', columns = ``to'')}
\]

\textbf{Step 2 --- Set up steady-state equation $V(T - I) = 0$.}
\[
T - I = \begin{pmatrix} 0.8 & 0.2 \\ 0.6 & 0.4 \end{pmatrix} - \begin{pmatrix} 1 & 0 \\ 0 & 1 \end{pmatrix} = \begin{pmatrix} -0.2 & 0.2 \\ 0.6 & -0.6 \end{pmatrix}
\]
\[
[S \;\; I]\begin{pmatrix} -0.2 & 0.2 \\ 0.6 & -0.6 \end{pmatrix} = [0 \;\; 0]
\]

\textbf{Step 3 --- Extract equation from column 1.}
\[
-0.2S + 0.6I = 0 \;\Rightarrow\; S = 3I \quad \cdots (1)
\]
(Column 2 gives $0.2S - 0.6I = 0$ --- same equation. Discard.)

\textbf{Step 4 --- Apply normalization constraint $S + I = 1$.}
\[
S + I = 1 \quad \cdots (2)
\]

\textbf{Step 5 --- Solve.}

Substituting (1) into (2):
\[
3I + I = 1 \;\Rightarrow\; 4I = 1 \;\Rightarrow\; \boxed{I = 0.25}
\]
\[
S = 3I = 3 \times 0.25 \;\Rightarrow\; \boxed{S = 0.75}
\]

\textbf{Steady-state vector:}
\[
V = [S \;\; I] = [0.75 \;\; 0.25]
\]

\textbf{Interpretation:} In the long run, the system is in the \textbf{Secure} state \textbf{75\%} of the time and in the \textbf{Insecure} state only \textbf{25\%} of the time.

\[
\boxed{\text{The system is more likely SECURE in the long run (75\% vs 25\%).}}
\]

\textit{Intuition: even when insecure, the 60\% patch-recovery rate is strong. Combined with 80\% staying-secure rate, the patching mechanism is effective overall --- the system gravitates toward the secure state.}
\end{solbox}

%% ==========================================
\section*{2020--2022 --- No Markov Chain question found}
%% ==========================================

\begin{warnbox}{2020/2021/2022 papers}
No Markov Chain question was found in the 2020, 2021, or 2022 past papers for CSE-717 / CSE-7197. The topic first appears in 2023 and 2024 (consecutive most recent years) --- a strong signal that it has been added to the repeating question pool and will very likely appear in 2026.
\end{warnbox}

%% ==========================================
\section*{Exam Strategy Summary}
%% ==========================================

\begin{keybox}{High-Yield Checklist for Markov Chains (2/5 years, both most recent, 3--4 marks)}
\begin{enumerate}[noitemsep]
  \item \textbf{Transition matrix orientation:} rows = ``from'' state, columns = ``to'' state. Each \textbf{row sums to 1}. Check this immediately after writing the matrix.
  \item \textbf{Fill in the missing probability:} if ``Secure$\to$Secure $= 0.4$'' is given, then ``Secure$\to$Compromised $= 0.6$'' automatically. Never leave a row incomplete.
  \item \textbf{Steady-state equation:} $VT = V \;\Leftrightarrow\; V(T-I) = 0$. Expand to get one useful linear equation (the two equations are always dependent --- keep only one).
  \item \textbf{Normalization:} $\pi_1 + \pi_2 = 1$ is the second equation. Together with the linear relation, solve in two lines.
  \item \textbf{Always interpret the result:} state which system (Secure vs Compromised/Insecure) dominates, with the exact long-run percentage. This interpretation is the final mark.
  \item \textbf{Both exam questions are 2-state security models} (Secure/Compromised in 2024, Secure/Insecure in 2023). Expect the same framing in 2026 --- the lecture worked example is a weather chain (Sunny/Rainy), but map the method directly onto the security context.
\end{enumerate}

\medskip
\textbf{Key contrast between the two years:}
\begin{center}
\begin{tabular}{lcc}
\toprule
 & \textbf{2024} & \textbf{2023} \\
\midrule
States & Secure / Compromised & Secure / Insecure \\
Stay-in-Secure prob & 0.4 & 0.8 \\
Recover-to-Secure prob & 0.3 & 0.6 \\
$\pi_{\text{Secure}}$ & $1/3 \approx 0.333$ & $3/4 = 0.75$ \\
$\pi_{\text{Compromised/Insecure}}$ & $2/3 \approx 0.667$ & $1/4 = 0.25$ \\
Long-run verdict & Compromised & Secure \\
\bottomrule
\end{tabular}
\end{center}
\end{keybox}

\end{document}
